% =====================================================================
%  BÖLÜM 8: Merkezi Limit Teoremi ve Asimptotik Teori
%  Kaynak: Feller Vol.II; Billingsley; Durrett; Klenke; Gnedenko
% =====================================================================
\chapter{Merkezi Limit Teoremi ve Asimptotik Teori}
\label{chp:merkezi_limit}

\bolumalinti{%
  Merkezi Limit Teoremi, olasılık teorisinin en güzel ve en derin sonuçlarından biridir.
  Normal dağılım, sanki evrenin derinliklerine işlenmiş gibidir.%
}{William Feller, \textit{An Introduction to Probability Theory and Its Applications}, Cilt II}

% ---- Kazanımlar (TYMM) ----
\begin{kazanimlar}
Bu bölümün sonunda öğrenci:
\begin{itemize}[leftmargin=1.8em]
  \item Merkezi Limit Teoremi'ni (MLT) ifade edip karakteristik fonksiyon yöntemiyle ispatlayabilir;
        De Moivre-Laplace özel durumunu tarihsel bağlamda açıklayabilir.
  \item Lindeberg koşulunu tanımlayabilir; bağımsız ama özdeş dağılımlı olmayan
        değişkenler için MLT'nin geçerli olduğu koşulları açıklayabilir.
  \item Çok boyutlu MLT'yi ifade edip kovaryans yapısıyla ilişkilendirebilir.
  \item Berry-Esseen teoremini kullanarak yakınsama hızına sayısal üst sınır verebilir.
  \item Yinelenmiş logaritma kanununun MLT ile ilişkisini sezgisel düzeyde açıklayabilir.
  \item Delta yöntemini büyük örneklem istatistiğine uygulayabilir;
        normal yaklaşım kalitesini değerlendirebilir.
\end{itemize}
\end{kazanimlar}

% ---- Motivasyon ----
\begin{motivasyon}
Büyük Sayılar Kanunu der ki: $\bar{X}_n \to \mu$. Soru şu: bu yakınsama \emph{ne kadar hızlı}
gerçekleşir ve $\bar{X}_n$ ile $\mu$ arasındaki fark \emph{nasıl dağılır}?

Merkezi Limit Teoremi yanıtlar: $\sqrt{n}(\bar{X}_n - \mu)/\sigma$ dağılımı,
$X_i$'lerin dağılımından \emph{bağımsız} olarak Normal(0,1)'e yakınsar.
Bu olağanüstü evrensellik, normal dağılımın olasılık ve istatistiğin merkezine
yerleşmesinin matematiksel gerekçesidir.

\medskip
\textbf{Köprü.} Bölüm~\ref{chp:beklenti}'deki karakteristik fonksiyon,
Bölüm~\ref{chp:donusumler}'deki delta yöntemi ve Bölüm~\ref{chp:buyuk_sayilar}'daki
Slutsky teoremi bu bölümün araçlarıdır.
Bölüm~\ref{chp:istatistik_temelleri}'nde büyük örneklem istatistiği bu sonuçlara dayanır.
\end{motivasyon}

% =====================================================================
\section{De Moivre-Laplace Teoremi}
\label{sec:demoivre_laplace}
% =====================================================================

\begin{tarihselnot}
Abraham de Moivre (1733) binom dağılımının büyük $n$ için normale yakınsadığını
fark etti. Pierre-Simon Laplace (1812) bunu daha geniş bir çerçevede ifade etti.
Carl Friedrich Gauss ise normal dağılımı hata teorisi bağlamında bağımsız olarak geliştirdi.
Genel i.i.d.\ versiyonu Aleksandr Lyapunov (1901) ve Jarl Waldemar Lindeberg (1922) tarafından
ispatlandı.
\end{tarihselnot}

\begin{teorem}[De Moivre-Laplace]
\label{thm:demoivre_laplace}
$X_i \sim \mathrm{Bernoulli}(p)$ bağımsız, $S_n = \sum_{i=1}^n X_i \sim \Binom(n,p)$.
$\mu = np$, $\sigma^2 = npq$. O zaman her $a < b$'de
\[
  P\!\left(a \leq \frac{S_n - np}{\sqrt{npq}} \leq b\right) \to \Phi(b) - \Phi(a) \quad (n\to\infty).
\]
\end{teorem}

\begin{proof}[İskelet İspat — Stirling Yaklaşımı]
Stirling formülü: $k! \approx \sqrt{2\pi k}(k/e)^k$.
$P(S_n = k)$ için $\binom{n}{k}p^kq^{n-k}$'yi Stirling ile açarak
$k = np + x\sqrt{npq}$ ($x$ sabit) için bu olasılığın
$\phi(x)/\sqrt{npq}$'ya yakınsadığı gösterilir (yerel limit teoremi).
Yerel limit toplamından dağılımla yakınsama elde edilir.
\end{proof}

\begin{ornek}[Binom Normal Yaklaşımı]
$n = 100$, $p = 0.4$: $\mu = 40$, $\sigma = \sqrt{24} \approx 4.90$.
$P(35 \leq S_{100} \leq 45) \approx \Phi\!\left(\frac{45.5-40}{4.90}\right) - \Phi\!\left(\frac{34.5-40}{4.90}\right)$
$= \Phi(1.12) - \Phi(-1.12) \approx 0.7372$.
(Süreklilik düzeltmesi: ayrık $\to$ sürekli, sınırlara $\pm 0.5$ eklenir.)
\end{ornek}

% =====================================================================
\section{Genel Merkezi Limit Teoremi}
\label{sec:genel_mlt}
% =====================================================================

\begin{teorem}[Merkezi Limit Teoremi — i.i.d.\ Durum]
\label{thm:mlt}
$X_1, X_2, \ldots$ i.i.d., $E[X_1] = \mu$, $\mathrm{Var}(X_1) = \sigma^2 \in (0,\infty)$. O zaman
\[
  Z_n = \frac{\bar{X}_n - \mu}{\sigma/\sqrt{n}} = \frac{\sum_{i=1}^n (X_i - \mu)}{\sigma\sqrt{n}}
  \xrightarrow{d} \Normal(0,1).
\]
\end{teorem}

\begin{proof}[Karakteristik Fonksiyon İspati]
$Y_i = (X_i - \mu)/\sigma$ tanımı: $E[Y_i] = 0$, $E[Y_i^2] = 1$.
$\varphi_Y(t) := E[e^{itY_1}]$ KF'si. Taylor açılımı ($E[Y_1]=0$, $E[Y_1^2]=1$):
\[
  \varphi_Y(t) = 1 - \frac{t^2}{2} + o(t^2) \quad (t\to 0).
\]
$Z_n = \frac{1}{\sqrt{n}}\sum_{i=1}^n Y_i$; bağımsızlık + KF çarpım:
\[
  \varphi_{Z_n}(t) = \left[\varphi_Y\!\left(\frac{t}{\sqrt{n}}\right)\right]^n
  = \left[1 - \frac{t^2}{2n} + o\!\left(\frac{1}{n}\right)\right]^n.
\]
$n\to\infty$ limiti: $\left(1 - \frac{t^2}{2n} + o(1/n)\right)^n \to e^{-t^2/2}$.
(Temel limit: $(1+a_n/n)^n \to e^a$ eğer $na_n \to a$.)
$e^{-t^2/2}$: standart normalin KF'si.
Lévy süreklilik teoremi (Teorem~\ref{prop:kf_yaklinsama}): $Z_n \xrightarrow{d} \Normal(0,1)$.
\end{proof}

\begin{kritiksorgu}
MLT hangi koşulları \emph{gerektirmez}?
\begin{itemize}
  \item $X_i$ normal olmak \emph{zorunda değil} (herhangi bir dağılım).
  \item $E[X_i^3] < \infty$ \emph{gerekmez} — yalnızca $\sigma^2 < \infty$ yeterli.
  \item $X_i$'lerin sürekli olması \emph{gerekmez} (kesikli dağılımlar da MLT'ye tabidir).
\end{itemize}
Ne gerekir: i.i.d. + sonlu varyans. Bunlar ortadan kalkarsa MLT bozulabilir
(Cauchy: varyans yok; ağır kuyruklar: stabil dağılımlara yakınsama).
\end{kritiksorgu}

\begin{onerme}[Yararlı Eşdeğer Formlar]
\label{prop:mlt_form}
MLT'den türetilen ifadeler:
\begin{align*}
  \sqrt{n}(\bar{X}_n - \mu) &\xrightarrow{d} \Normal(0, \sigma^2), \\
  \frac{S_n - n\mu}{\sigma\sqrt{n}} &\xrightarrow{d} \Normal(0,1), \\
  P\!\left(\bar{X}_n \leq \mu + \frac{z_\alpha \sigma}{\sqrt{n}}\right) &\to \Phi(z_\alpha) = 1-\alpha.
\end{align*}
\end{onerme}

% =====================================================================
\section{Lindeberg-Feller Koşulları}
\label{sec:lindeberg}
% =====================================================================

\begin{kesif}
i.i.d.\ koşulunu kaldırırsak ne olur? $X_i$'lerin farklı dağılımları olsa bile
toplamları normale yakınsar mı? Hangi koşul "hiçbir terim baskın değil"
sezgisini matematiksel olarak yakalar?
\end{kesif}

\begin{tanim}[Lindeberg Koşulu]
\label{def:lindeberg}
$\{X_{ni}\}$ bağımsız, $E[X_{ni}] = 0$, $\sigma_{ni}^2 = \mathrm{Var}(X_{ni})$,
$s_n^2 = \sum_{i=1}^n \sigma_{ni}^2$. \textbf{Lindeberg koşulu}: her $\varepsilon > 0$ için
\[
  L_n(\varepsilon) = \frac{1}{s_n^2} \sum_{i=1}^n E\!\left[X_{ni}^2 \mathbf{1}_{|X_{ni}|>\varepsilon s_n}\right] \to 0.
\]
\end{tanim}

\begin{teorem}[Lindeberg-Feller MLT]
\label{thm:lindeberg_feller}
Lindeberg koşulu sağlanıyorsa
$\frac{1}{s_n}\sum_{i=1}^n X_{ni} \xrightarrow{d} \Normal(0,1)$.
Tersine (Feller): Lindeberg koşulu, $\max_i \sigma_{ni}/s_n \to 0$ ("ufaklık koşulu")
altında MLT için zorunludur.
\end{teorem}

\begin{proof}[Ana Fikir]
Lindeberg koşulu, her terimin katkısının $s_n$'e kıyasla küçük olmasını garantiler.
KF ispatında $\varphi_{X_{ni}/s_n}(t) = 1 - t^2\sigma_{ni}^2/(2s_n^2) + R_{ni}(t)$;
hata terimi $|R_{ni}(t)| \leq E[X_{ni}^2\mathbf{1}_{|X_{ni}|>\varepsilon s_n}] \cdot t^2/s_n^2$.
Lindeberg koşulu ile hataların toplamı sıfıra gider;
$\prod \varphi_{X_{ni}/s_n}(t) \to e^{-t^2/2}$ sonucuna varılır.
Ayrıntı: Billingsley~\cite{Billingsley1995}, §27.
\end{proof}

\begin{onerme}[Lyapunov Koşulu $\Rightarrow$ Lindeberg]
\label{prop:lyapunov}
$\sum_{i=1}^n E[|X_{ni}|^{2+\delta}] / s_n^{2+\delta} \to 0$ (bazı $\delta > 0$)
$\Rightarrow$ Lindeberg koşulu $\Rightarrow$ MLT.
Lyapunov koşulu daha hesaplanabilir ama Lindeberg'den daha kuvvetlidir.
\end{onerme}

\begin{ornek}[Lindeberg Koşulunun Önemi]
$n$ işçinin maaş artışı: $X_i \sim \Normal(0, \sigma_i^2)$, $\sigma_i^2 = i$.
$s_n^2 = \sum i = n(n+1)/2$.
Lyapunov ($\delta=1$): $\sum E[|X_i|^3]/s_n^3 \propto \sum i^{3/2}/n^{3/2} \to 0$.
Toplam büyük olduğunda bile MLT geçerli; Lindeberg sağlanır.
\end{ornek}

% =====================================================================
\section{Çok Boyutlu Merkezi Limit Teoremi}
\label{sec:cok_boyut_mlt}
% =====================================================================

\begin{teorem}[Çok Boyutlu MLT]
\label{thm:cok_boyut_mlt}
$\mathbf{X}_1, \mathbf{X}_2, \ldots$ i.i.d.\ $\mathbb{R}^k$ değerli,
$E[\mathbf{X}_i] = \boldsymbol{\mu}$, kovaryans matrisi $\boldsymbol{\Sigma}$ (pozitif yarı-tanımlı). O zaman
\[
  \sqrt{n}(\bar{\mathbf{X}}_n - \boldsymbol{\mu}) \xrightarrow{d} \Normal_k(\mathbf{0},\, \boldsymbol{\Sigma}).
\]
\end{teorem}

\begin{proof}
\textbf{Cramér-Wold cihazı:} $\mathbf{Y}_n \xrightarrow{d} \mathbf{Y}$ ($\mathbb{R}^k$'de) ancak ve ancak
her $\mathbf{a} \in \mathbb{R}^k$ için $\mathbf{a}^\top \mathbf{Y}_n \xrightarrow{d} \mathbf{a}^\top \mathbf{Y}$.

$\mathbf{a}^\top \sqrt{n}(\bar{\mathbf{X}}_n - \boldsymbol{\mu}) = \sqrt{n}(\overline{\mathbf{a}^\top\mathbf{X}_n} - \mathbf{a}^\top\boldsymbol{\mu})$.

$\mathbf{a}^\top\mathbf{X}_i$: i.i.d., ortalama $0$, varyans $\mathbf{a}^\top\boldsymbol{\Sigma}\mathbf{a}$.
Tek boyutlu MLT: $\xrightarrow{d} \Normal(0, \mathbf{a}^\top\boldsymbol{\Sigma}\mathbf{a})$.
Bu $\Normal_k(\mathbf{0},\boldsymbol{\Sigma})$'nın doğrusal kombinasyonudur; Cramér-Wold tamamlar.
\end{proof}

\begin{ornek}[Çok Boyutlu Normal Yaklaşım]
$(X_i, Y_i)$ i.i.d., $E[\mathbf{X}] = (\mu_X, \mu_Y)^\top$,
$\boldsymbol{\Sigma} = \begin{pmatrix}\sigma_X^2 & \rho\sigma_X\sigma_Y \\ \rho\sigma_X\sigma_Y & \sigma_Y^2\end{pmatrix}$.
Büyük $n$: $(\bar{X}_n, \bar{Y}_n)$ yaklaşık $\Normal_2\!\left((\mu_X,\mu_Y)^\top,\, \boldsymbol{\Sigma}/n\right)$.
Bu, iki boyutlu güven bölgelerinin (elips) temelidir.
\end{ornek}

% =====================================================================
\section{Berry-Esseen Teoremi}
\label{sec:berry_esseen}
% =====================================================================

\begin{motivasyon}[MLT'nin Hızı Nedir?]
MLT $Z_n \xrightarrow{d} \Normal(0,1)$ diyor. Ama $n = 30$ iken ne kadar yakınız?
Berry-Esseen teoremi hata büyüklüğüne sayısal üst sınır koyar.
\end{motivasyon}

\begin{teorem}[Berry-Esseen]
\label{thm:berry_esseen}
$X_1, \ldots, X_n$ i.i.d., $E[X_1] = 0$, $E[X_1^2] = \sigma^2 > 0$,
$E[|X_1|^3] = \rho < \infty$. O zaman her $x \in \mathbb{R}$ için
\[
  \left|F_{Z_n}(x) - \Phi(x)\right| \leq \frac{C\,\rho}{\sigma^3\sqrt{n}},
\]
burada $C \leq 0.4748$ (en iyi bilinen sabit; Shevtsova 2010).
\end{teorem}

İspatın ana fikri: KF'ler farkına Fourier ters dönüşümü uygulanması;
"yumuşatma" (smoothing) lemması aracılığıyla KDF farkına bağlanması.
Ayrıntı için Feller Vol.II~\cite{Feller1971}, Bölüm XVI.5 veya Durrett~\cite{Durrett2019}, §3.4.

\begin{ornek}[Berry-Esseen Uygulaması]
$X_i \sim \mathrm{Bernoulli}(0.5)$: $\sigma^2 = 0.25$, $\rho = E[|X_1|^3] = 0.5$.
$n = 36$: hata $\leq C \cdot 0.5/(0.25^{3/2}\cdot 6) = C \cdot 0.5/(0.125\cdot 6) \approx 0.663C \approx 0.315$.
$n = 100$: hata $\leq C\cdot 0.5/(0.125\cdot 10) \approx 0.190$.
$n = 1000$: hata $\leq 0.060$: 3 ondalık basamak için $n \approx 1000$ önerilir.
\end{ornek}

\begin{hataanalizi}
"$n \geq 30$ ise MLT yeterince iyi yaklaşır" — çok kaba bir kuraldır.
Berry-Esseen: hata üçüncü mutlak moment ile orantılı.
Çarpık dağılımlarda ($\rho$ büyük) $n = 30$ yetersizdir;
çok çarpık dağılımlar (lognormal, Pareto) için $n = 1000$ bile tartışmalıdır.
\end{hataanalizi}

% =====================================================================
\section{Poisson Yakınsaması}
\label{sec:poisson_yaklinsama}
% =====================================================================

\begin{teorem}[Le Cam Teoremi]
\label{thm:le_cam}
$X_1, \ldots, X_n$ \emph{bağımsız} (özdeş dağılımlı olmak zorunda değil),
$X_i \sim \mathrm{Bernoulli}(p_i)$, $S_n = \sum X_i$.
$\lambda_n = \sum p_i$. Toplam varyasyon mesafesi:
\[
  d_{TV}(S_n, \Pois(\lambda_n)) \leq 2\sum_{i=1}^n p_i^2 = 2\sum p_i^2.
\]
\end{teorem}

\begin{proof}[Ana Fikir — Stein-Chen Yöntemi]
Poisson operatörü: $\mathcal{A}f(k) = \lambda f(k+1) - kf(k)$.
$E[\mathcal{A}f(S_n)]$ ile $E[\mathcal{A}f(\Pois(\lambda))]$ arasındaki fark,
$\sum p_i^2$ terimlerinden beslenir.
Ayrıntı: Arratia-Goldstein-Gordon (1989) veya Durrett~\cite{Durrett2019}, §3.6.
\end{proof}

\begin{ornek}[Le Cam Uygulaması]
$n = 200$, $p_i = 1/200$ homojen: $\lambda = 1$, $\sum p_i^2 = 200 \cdot (1/200)^2 = 1/200$.
$d_{TV} \leq 2/200 = 0.01$: Poisson yaklaşımı $S_{200}$ için çok iyidir.
Bu, Bölüm~\ref{subsec:poisson}'daki binom-Poisson limitinin nicel versiyonudur.
\end{ornek}

% =====================================================================
\section{Yinelenmiş Logaritma Kanunu}
\label{sec:yinelenmis_log}
% =====================================================================

\begin{motivasyon}[MLT'nin Ötesi]
MLT: $S_n/({\sigma\sqrt{n}}) \xrightarrow{d} \Normal(0,1)$. Yani $S_n \approx \sigma\sqrt{n}$ büyüklüğünde.
Ama $S_n/(\sigma\sqrt{n})$'nin üstel rastlanımsal salınımı ne kadar büyük olabilir?
Sezgi: bazen $k\sigma\sqrt{n}$'ye ($k > 1$) çıkabilir; bu "bazen"in kesin sınırı nedir?
\end{motivasyon}

\begin{teorem}[Yinelenmiş Logaritma Kanunu (YLK)]
\label{thm:ylk}
$X_1, X_2, \ldots$ i.i.d., $E[X_1] = 0$, $\mathrm{Var}(X_1) = \sigma^2 \in (0,\infty)$,
$S_n = \sum_{i=1}^n X_i$. O zaman
\[
  \limsup_{n\to\infty} \frac{S_n}{\sigma\sqrt{2n\ln\ln n}} = 1 \quad \text{n.k.}
\]
ve $\liminf_{n\to\infty} S_n/(\sigma\sqrt{2n\ln\ln n}) = -1$ n.k.
\end{teorem}

İspat oldukça teknik olup Kolmogorov'un orijinal 1929 kanıtına, daha sonra
Hartman-Wintner (1941) genel versiyonuna dayanır.
Temel fikir: uygun olayların bağımsızlığı + Borel-Cantelli lemmalarının dikkatli uygulanması.

\begin{onerme}[YLK ve MLT Karşılaştırması]
\begin{center}
\begin{tabular}{lcc}
\toprule
 & MLT & YLK \\
\midrule
Normalize eden & $\sigma\sqrt{n}$ & $\sigma\sqrt{2n\ln\ln n}$ \\
Limit türü & Dağılımla ($\Normal(0,1)$) & n.k.\ ($\limsup = 1$) \\
Anlam & Tipik davranış & Ekstrem salınım \\
\bottomrule
\end{tabular}
\end{center}
\end{onerme}

\begin{ornek}[YLK Yorumu]
$X_i \sim \Normal(0,1)$ bağımsız. $S_n$ neredeyse kesinlikle
$\pm\sqrt{2n\ln\ln n}$ bandının içinde kalır; bu bandı sonsuz kez tam olarak dokunur
ama aşmaz (kesin üst sınır).
$n = 10^6$: bant genişliği $\approx \sqrt{2 \times 10^6 \times \ln\ln(10^6)} \approx \sqrt{2\times10^6\times2.93} \approx 2420$.
MLT bandı $\sqrt{n} \approx 1000$; YLK bant $\approx 2.4$ kat daha geniş.
\end{ornek}

% =====================================================================
\section{Büyük Örneklemde Normal Yaklaşımlar}
\label{sec:buyuk_orneklem}
% =====================================================================

\begin{onerme}[Örneklem Ortalaması — Güven Aralığı]
\label{prop:guven_araligi}
$X_1, \ldots, X_n$ i.i.d., $\mu$ bilinmiyor, $\sigma^2$ biliniyor. MLT'den:
\[
  P\!\left(\bar{X}_n - z_{\alpha/2}\frac{\sigma}{\sqrt{n}} \leq \mu \leq \bar{X}_n + z_{\alpha/2}\frac{\sigma}{\sqrt{n}}\right) \to 1-\alpha.
\]
$z_{0.025} = 1.96$: yaklaşık $\%95$ güven aralığı $\bar{X}_n \pm 1.96\sigma/\sqrt{n}$.
\end{onerme}

\begin{onerme}[Örneklem Oranı]
\label{prop:oran_guven}
$X_i \sim \mathrm{Bernoulli}(p)$, $\hat{p} = \bar{X}_n$. MLT:
$\sqrt{n}(\hat{p}-p)/\sqrt{p(1-p)} \xrightarrow{d} \Normal(0,1)$.
$p$ bilinmiyorsa $\sqrt{\hat{p}(1-\hat{p})}$ ile değiştirin (Slutsky):
yaklaşık $\%95$ GA: $\hat{p} \pm 1.96\sqrt{\hat{p}(1-\hat{p})/n}$.
\end{onerme}

\begin{proof}
$\hat{p}$ yerine $p$ yazılan versiyonu MLT'nin doğrudan uygulaması.
$\sqrt{\hat{p}(1-\hat{p})} \xrightarrow{P} \sqrt{p(1-p)}$ (BSK + sürekli haritalama).
Slutsky: bölüm dağılımla $\Normal(0,1)$'e yakınsar.
\end{proof}

\begin{teorem}[Delta Yöntemi — MLT Uygulaması]
\label{thm:delta_mlt}
$\sqrt{n}(\bar{X}_n - \mu) \xrightarrow{d} \Normal(0,\sigma^2)$ ve $g$, $\mu$'da türevlenebilir ise
(Teorem~\ref{thm:delta_yontemi} ile birlikte):
\[
  \sqrt{n}(g(\bar{X}_n) - g(\mu)) \xrightarrow{d} \Normal(0,\, \sigma^2[g'(\mu)]^2).
\]
Bu, Bölüm~\ref{chp:nokta_tahmin}'deki asimptotik güven aralıklarının temelidir.
\end{teorem}

\begin{ornek}[Varyans Stabilize Eden Dönüşüm]
$X_i \sim \mathrm{Bernoulli}(p)$: $\sigma^2 = p(1-p)$, $p$'ye bağlı.
$g(p) = \arcsin(\sqrt{p})$: $g'(p) = 1/(2\sqrt{p(1-p)})$.
Delta: $[g'(p)]^2 \sigma^2 = p(1-p)/(4p(1-p)) = 1/4$ — sabit!
$\sqrt{n}(\arcsin\sqrt{\hat{p}} - \arcsin\sqrt{p}) \xrightarrow{d} \Normal(0, 1/4)$.
Varyans artık $p$'den bağımsız: güven aralıkları her $p$ için eşit genişlikte.
\end{ornek}

\begin{onerme}[Poisson Gözlemleri için Normal Yaklaşım]
$X \sim \Pois(\lambda)$, $\lambda$ büyük:
$(\sqrt{X} - \sqrt{\lambda}) \xrightarrow{d} \Normal(0, 1/4)$ (karekök dönüşümü).
$P(X \leq k) \approx \Phi\!\left(\frac{k+0.5-\lambda}{\sqrt{\lambda}}\right)$ (süreklilik düzeltmesi).
\end{onerme}

% ---- Disiplinlerarası Bağlantı (TYMM) ----
\begin{kopru}[Disiplinlerarası — Fizik, Sosyal Bilimler, Biyoistatistik]
\textbf{Fizik:} Termal dengedeki gaz moleküllerinin Maxwell-Boltzmann hız dağılımı,
çok sayıda bağımsız çarpışmanın MLT sonucu olarak Gaussian'a yakınsar.
Hata teorisi (Gauss): ölçüm hatalarının bağımsız küçük bileşenlerden oluştuğu
varsayımı altında toplam hata normal dağılır.\\
\textbf{Sosyal bilimler:} Büyük örneklem istatistiği (anket, seçim tahminleri, A/B testi)
%\%95 güven aralığı MLT'ye dayanır. "Hata payı $\pm 3\%$" beyanı normal yaklaşımın uygulamasıdır.\\
\textbf{Biyoistatistik:} Klinik denemelerde $t$-testi, MLT'nin küçük örneklemlerde
doğrudan uygulaması değil, $t$ dağılımının kullanımıdır; büyük $n$'de $t_n \to \Normal(0,1)$.
\end{kopru}

% ---- Öz-Yansıtma ----
\begin{ozyansima}
\begin{itemize}[leftmargin=1.8em]
  \item MLT'nin "evrenselliği" size sezgisel olarak neden doğru geliyor (ya da gelmiyor)?
  \item Berry-Esseen teoremi "$n \geq 30$ yeterli" kuralını nasıl sorgular?
  \item Lindeberg koşulu intuisyon olarak ne anlama gelir?
        "Hiçbir terim baskın değil" yorumunu bir örnekle somutlaştırın.
  \item YLK ile MLT arasındaki farkı bir grafikle anlatın (sözlü olarak tarif edin).
\end{itemize}
\end{ozyansima}

% ---- Köprü Notu ----
\begin{kopru}
\textbf{Önceki bölüm:} Bölüm~\ref{chp:buyuk_sayilar} BSK'yı verdi (yakınsama garantisi);
bu bölüm yakınsama \emph{hızını} ve \emph{dağılımını} belirledi.\\
\textbf{Sonraki (opsiyonel):} Bölüm~\ref{chp:martingaller} — koşullu beklenti
ve martingallar; martingal MLT.\\
\textbf{İstatistik kısmı:} Bölüm~\ref{chp:istatistik_temelleri}'nde asimptotik
dağılım teorisi; Bölüm~\ref{chp:nokta_tahmin}'de MLE'nin asimptotik normalliği
($\sqrt{n}(\hat\theta_{MLE}-\theta) \xrightarrow{d} \Normal(0, I(\theta)^{-1})$).
\defterbaglantisi{bolum08\_mlt.ipynb}{%
  MLT simülasyonu (çeşitli dağılımlar), Berry-Esseen görselleştirmesi, YLK yolu}
\end{kopru}

% =====================================================================
%  ALIŞTIRMALAR
% =====================================================================

\cozumlualistirmalar

\begin{alistirma}[Al.8.1]{A}{MLT Uygulaması}
Bir sigorta şirketinin yıllık hasar başına ödeme miktarı $E[X] = 1000$ TL,
$\sigma = 200$ TL. $n = 100$ hasar için toplam ödemenin $105\,000$ TL'yi aşma
olasılığını MLT ile tahmin edin.
\end{alistirma}
\alcozum{%
$S_{100}$: $E[S_{100}] = 100\,000$, $\mathrm{Var}(S_{100}) = 100 \times 200^2 = 4\,000\,000$,
$\sigma_{S_{100}} = 2000$.
$P(S_{100} > 105\,000) = P\!\left(Z > \frac{105000-100000}{2000}\right) = P(Z > 2.5) = 1-\Phi(2.5) \approx 0.0062$.}

\begin{alistirma}[Al.8.2]{A}{De Moivre-Laplace}
$S_{400} \sim \Binom(400, 0.4)$.
(a) $P(150 \leq S_{400} \leq 170)$'i MLT ile hesaplayın (süreklilik düzeltmesiyle).
(b) $E[S_{400}]$ ve $\sigma_{S_{400}}$'ü hesaplayın.
\end{alistirma}
\alcozum{%
(b) $E[S_{400}]=160$, $\sigma=\sqrt{400\cdot0.4\cdot0.6}=\sqrt{96}\approx9.798$.\\
(a) Süreklilik düzeltmesiyle: $P(149.5\leq S\leq170.5)$.
$z_1=(149.5-160)/9.798\approx-1.072$, $z_2=(170.5-160)/9.798\approx1.072$.
$P\approx\Phi(1.072)-\Phi(-1.072)=2\Phi(1.072)-1\approx2(0.8582)-1=0.7164$.}

\begin{alistirma}[Al.8.3]{B}{Berry-Esseen Sınırı}
$X_i \sim \mathrm{Exp}(1)$ bağımsız: $\mu=1$, $\sigma^2=1$, $E[|X_1-1|^3]$'ü hesaplayın.
Berry-Esseen ile $n=50$ için maksimum KDF hatasına üst sınır verin.
\end{alistirma}
\alcozum{%
$Y_i = X_i-1$ için $E[|Y_i|^3]$.
$E[|X-1|^3] = \int_0^\infty |x-1|^3 e^{-x}dx = \int_0^1(1-x)^3e^{-x}dx + \int_1^\infty(x-1)^3e^{-x}dx$.
Entegrasyon by parts (veya mgf türevi): $E[|X-1|^3] = 2$ (Exp(1) için).
$\sigma^3 = 1$, $\rho = 2$.
Berry-Esseen: hata $\leq C\cdot2/(\sqrt{50})\approx 0.4748\cdot2/7.071\approx0.134$.
$n=50$, üstel dağılım için maksimum KDF hatası yaklaşık $0.134$.}

\begin{alistirma}[Al.8.4]{B}{Delta Yöntemi + MLT}
$X_1,\ldots,X_n$ i.i.d.\ $\mathrm{Pois}(\lambda)$, $\bar{X}_n$ örneklem ortalaması.
(a) MLT: $\sqrt{n}(\bar{X}_n-\lambda)$'nın limit dağılımını yazın.
(b) $g(x) = \sqrt{x}$ dönüşümü ile $\sqrt{n}(\sqrt{\bar{X}_n}-\sqrt{\lambda})$'nın limit dağılımını bulun.
(c) (b)'nin sonucu neden uygulamada faydalıdır?
\end{alistirma}
\alcozum{%
(a) $\mathrm{Var}(X_i)=\lambda$: $\sqrt{n}(\bar{X}_n-\lambda)\xrightarrow{d}\Normal(0,\lambda)$.\\
(b) $g'(x)=1/(2\sqrt{x})$, $g'(\lambda)=1/(2\sqrt{\lambda})$.
$\sqrt{n}(\sqrt{\bar{X}_n}-\sqrt{\lambda})\xrightarrow{d}\Normal(0,\lambda/(4\lambda))=\Normal(0,1/4)$.\\
(c) Varyans $1/4$ artık $\lambda$'ya bağlı değil: güven aralığı $\sqrt{\hat\lambda}\pm z_{\alpha/2}/(2\sqrt{n})$.
Bilinmeyen parametre ile güven aralığı genişliğini kontrol etmek mümkün.}

% ---

\alistirmalar

\begin{alistirma}[Al.8.5]{A}{Örneklem Oranı}
$n=400$ kişilik ankette $180$ kişi "evet" dedi. $\hat{p}=0.45$.
(a) $p$ için yaklaşık $\%95$ güven aralığını hesaplayın.
(b) $p=0.5$ hipotezini $\alpha=0.05$ için test edin ($z$-testi).
\end{alistirma}
\alcozum{%
(a) $\hat{p}=0.45$, $\sqrt{\hat{p}(1-\hat{p})/n}=\sqrt{0.2475/400}=\sqrt{0.000619}\approx0.02488$.
GA: $0.45\pm1.96\times0.02488 = (0.401, 0.499)$.
(b) $z=(0.45-0.5)/\sqrt{0.5\times0.5/400}=(-0.05)/0.025=-2.0$.
$|z|=2.0>1.96$: $H_0$ reddedilir, $\alpha=0.05$.}

\begin{alistirma}[Al.8.6]{A}{Poisson Normal Yaklaşımı}
$X \sim \Pois(25)$. MLT ile $P(X \leq 30)$ için normal yaklaşım verin.
Süreklilik düzeltmesini kullanın.
\end{alistirma}
\alcozum{%
$E[X]=25$, $\sigma=5$.
$P(X\leq30)\approx P\!\left(Z\leq\frac{30.5-25}{5}\right)=\Phi(1.10)\approx0.8643$.
(Gerçek değer $\approx0.870$; süreklilik düzeltmesi olmadan $\Phi(1.0)\approx0.841$, daha az doğru.)}

\begin{alistirma}[Al.8.7]{B}{Lindeberg Koşulu}
$X_{ni} = \pm 1/\sqrt{n}$ eşit olasılıkla ($i = 1,\ldots,n$), bağımsız.
(a) $E[X_{ni}]=0$, $\sigma_{ni}^2=1/n$, $s_n^2=1$ olduğunu gösterin.
(b) Lindeberg koşulunun sağlandığını kontrol edin.
(c) $\sum X_{ni}$'nin limit dağılımını belirtin.
\end{alistirma}
\alcozum{%
(a) $\sigma_{ni}^2 = E[X_{ni}^2]=1/n$; $s_n^2=n\cdot(1/n)=1$. \\
(b) $L_n(\varepsilon) = \sum_i E[X_{ni}^2\mathbf{1}_{|X_{ni}|>\varepsilon s_n}]/s_n^2$.
$|X_{ni}|=1/\sqrt{n}$; $\varepsilon s_n=\varepsilon$.
$n$'yeterince büyüyünce $1/\sqrt{n}<\varepsilon$: indikatör sıfır.
$L_n(\varepsilon)=0$ (tüm büyük $n$). Sağlandı.\\
(c) Lindeberg-Feller: $\sum X_{ni}\xrightarrow{d}\Normal(0,1)$.}

\begin{alistirma}[Al.8.8]{B}{Çok Boyutlu MLT}
$(X_i, Y_i)$ i.i.d., $E[X]=E[Y]=0$, $\mathrm{Var}(X)=\mathrm{Var}(Y)=1$, $\mathrm{Cov}(X,Y)=\rho$.
$(\bar{X}_n, \bar{Y}_n)$'nin büyük örneklem dağılımını yazın.
$U_n = \bar{X}_n + \bar{Y}_n$'nin yaklaşık dağılımını bulun.
\end{alistirma}
\alcozum{%
$\boldsymbol{\Sigma} = \begin{pmatrix}1&\rho\\\rho&1\end{pmatrix}$.
$\sqrt{n}(\bar{X}_n,\bar{Y}_n)^\top \xrightarrow{d} \Normal_2(\mathbf{0},\boldsymbol{\Sigma})$.\\
$U_n = \bar{X}_n+\bar{Y}_n$: $\mathbf{a}=(1,1)^\top$.
$\mathrm{Var}(U_n)=\mathbf{a}^\top\boldsymbol{\Sigma}\mathbf{a}/n=(1+1+2\rho)/n=2(1+\rho)/n$.
$\sqrt{n}\,U_n\xrightarrow{d}\Normal(0,2(1+\rho))$.}

\begin{alistirma}[Al.8.9]{C}{MLT İspat Detayı}
$(1-t^2/(2n))^n \to e^{-t^2/2}$ limitini $n\to\infty$ için kesin biçimde ispatlayın.
Daha genel olarak: $a_n \to 0$ ve $na_n \to a$ iken $(1+a_n)^n \to e^a$ olduğunu gösterin.
\end{alistirma}
\alcozum{%
$\ln(1+a_n)^n = n\ln(1+a_n)$.
$\ln(1+x) = x - x^2/2 + O(x^3)$ ($x\to0$).
$n\ln(1+a_n) = n(a_n - a_n^2/2 + O(a_n^3))$.
$na_n \to a$; $na_n^2 = (na_n)\cdot a_n \to a\cdot0 = 0$ ($a_n\to0$).
$nO(a_n^3) = na_n \cdot a_n^2 \to 0$.
$\Rightarrow n\ln(1+a_n) \to a$.
$\Rightarrow (1+a_n)^n = e^{n\ln(1+a_n)} \to e^a$.
Özel durum: $a_n = -t^2/(2n)$, $na_n = -t^2/2 \to -t^2/2$.
$(1-t^2/(2n))^n \to e^{-t^2/2}$. \checkmark}

% =====================================================================
\endinput
